% Part: proof-theory
% Chapter: natural-deduction
% Section: quantifiers

\documentclass[../../../include/open-logic-section]{subfiles}

\begin{document}

\olfileid{pt}{nat}{qua}

\olsection{انتظام \usetoken{P}{proof} والتعويض}

تُحدث قواعد الأسوار بعض التعقيدات بسبب «شروط المتغير المميز» التي تنطبق
على~$\Elim{\lexists}$ و$\Intro{\lforall}$. ويمكن معالجة معظم هذه
التعقيدات بضمان عدم إعادة استخدام~!!{constant}~$c$ في هذه
الاستدلالات. فلنقدم التعريف الآتي.

\begin{defn}
يكون~!!{proof} الاستنباط الطبيعي~$\delta$ \emph{منتظمًا} إذا
\begin{enumerate}
  \item لم يُستخدم أي متغير مميز~$a$ متغيرًا مميزًا في أكثر من
  استدلال واحد لـ$\Elim{\lexists}$ أو~$\Intro{\lforall}$، و
  \item إذا كان~$a$ المتغير المميز في استدلال~$\Elim{\lexists}$ أو
  $\Intro{\lforall}$، لم يظهر إلا في~!!{proof} الجزئي المباشر المؤدي
  إلى المقدمة الصغرى أو إلى المقدمة، على الترتيب.
\end{enumerate}
\end{defn}

\begin{lem}\ollabel{lem:subst-regular} إذا كان~$\delta$ منتظمًا،
وينتهي بـ$!C$، وكان~$c$~!!a{constant} لا يُستخدم متغيرًا مميزًا في
$\delta$، وكان~$t$ حدًا لا يحتوي أي متغير مميز من~$\delta$، فإن
$\Subst{\delta}{t}{c}$ الناتج من~$\delta$ باستبدال كل ظهور لـ$c$ بـ
$t$ هو~!!{proof} منتظم. وتكون~!!{formula} نهاية
$\Subst{\delta}{t}{c}$ هي~$\Subst{!C}{t}{c}$. وإذا كان~$!A^x$ فرضًا
مفتوحًا في~$\delta$، كان~$\Subst{!A}{t}{c}^x$ فرضًا مفتوحًا في
$\Subst{\delta}{t}{c}$.
\end{lem}

\begin{proof}
  بالاستقراء على~$\pheight{\delta}$. إذا كان~$\pheight{\delta}=0$، لم
  يتألف إلا من الفرض~$!A^x$. وعندئذ يكون~$\Subst{\delta}{t}{c}$ هو
  $\Subst{!A}{t}{c}^x$.

  وإذا كان~$\pheight{\delta}>0$، ميزنا الحالات بحسب آخر استدلال في
  $\delta$. وحالات قواعد الروابط القضوية سهلة فنتركها تمارين. ونتأمل
  الحالات التي ينتهي فيها~$\delta$ باستدلال سُوَري.
  \begin{enumerate}
    \item إذا كان الاستدلال الأخير~$\Intro{\lforall}$، كان~$\delta$
    على الصورة
    \[
    \AxiomC{}
    \RightLabel{$\delta'$}
    \DeduceC{$!A(a,c)$}
    \RightLabel{\Intro{\lforall}}
    \UnaryInfC{$\lforall[x][!A(x,c)]$}
    \DisplayProof
    \]
    وبفرض الاستقراء يكون~$\Subst{\delta'}{t}{c}$~!!{proof} منتظمًا لـ
    $!A(a,t)$. ولأن~$a$ متغير مميز في~$\delta$، فلا تظهر~$a$ في~$t$.
    ومن ثم يكون الاستدلال الأخير في~$\Subst{\delta}{t}{c}$،
    \[
    \AxiomC{}
    \RightLabel{$\Subst{\delta'}{t}{c}$}
    \DeduceC{$!A(a,t)$}
    \RightLabel{\Intro{\lforall}}
    \UnaryInfC{$\lforall[x][!A(x,t)]$}
    \DisplayProof
    \]
    صحيحًا: لما كان~$t$ لا يحتوي~$a$، لم تحتوها النتيجة ولا أي فرض
    مفتوح.
    \item إذا كان الاستدلال الأخير~$\Elim{\lforall}$، كان~$\delta$
    على الصورة
    \[
    \AxiomC{}
    \RightLabel{$\delta'$}
    \DeduceC{$\lforall[x][!A(x,c)]$}
    \RightLabel{\Elim{\lforall}}
    \UnaryInfC{$!A(s(c),c)$}
    \DisplayProof
    \]
    وبفرض الاستقراء يكون~$\Subst{\delta'}{t}{c}$~!!{proof} منتظمًا لـ
    $\lforall[x][!A(x,t)]$. (وقد يحتوي الحد~$s(c)$ في النتيجة~$c$،
    وقد لا يحتويها.) ويكون الاستدلال الأخير في
    $\Subst{\delta}{t}{c}$،
    \[
    \AxiomC{}
    \RightLabel{$\Subst{\delta'}{t}{c}$}
    \DeduceC{$\lforall[x][!A(x,t)]$}
    \RightLabel{\Elim{\lforall}}
    \UnaryInfC{$!A(s(t),t)$}
    \DisplayProof
    \]
    صحيحًا، ويكون~$!A(s(t),t) \ident \Subst{!A(s(c),c)}{t}{c}$.
    وحالة~$\Intro{\lexists}$ مماثلة.
  \item إذا كان الاستدلال الأخير~$\Elim{\lexists}$، كان~$\delta$
    على الصورة
    \[
    \AxiomC{}
    \RightLabel{$\delta_1'$}
    \DeduceC{$\lexists[x][!A(x,c)]$}
    \AxiomC{$\Discharge{!A(a,c)}{x}$}
    \RightLabel{$\delta_2'$}
    \DeduceC{$!C$}
    \DischargeRule{\Elim{\lexists}}{x}
    \BinaryInfC{$!C$}
    \DisplayProof
    \]
    وبفرض الاستقراء يكون~$\Subst{\delta_1'}{t}{c}$ و
    $\Subst{\delta_2'}{t}{c}$~!!{proof}s منتظمين: الأول لـ
    $\lexists[x][!A(x,t)]$، والثاني لـ$\Subst{!C}{t}{c}$ من الفرض
    المفتوح~$!A(a,t)^x$. ولأن~$a$ متغير مميز في~$\delta$، فلا تظهر
    $a$ في~$t$. ومن ثم يكون الاستدلال الأخير في
    $\Subst{\delta}{t}{c}$،
    \[
    \AxiomC{}
    \RightLabel{$\Subst{\delta_1'}{t}{c}$}
    \DeduceC{$\lexists[x][!A(x,t)]$}
    \AxiomC{$\Discharge{!A(a,t)}{x}$}
    \RightLabel{$\Subst{\delta_2'}{t}{c}$}
    \DeduceC{$\Subst{!C}{t}{c}$}
    \DischargeRule{\Elim{\lexists}}{x}
    \BinaryInfC{$\Subst{!C}{t}{c}$}
    \DisplayProof
    \]
    صحيحًا: لا تظهر~$a$ في المقدمة الكبرى
    $\lexists[x][!A(x,t)]$، ولا في النتيجة~$\Subst{!C}{t}{c}$، ولا في
    أي فرض مفتوح سوى الفرض المبرأ~$!A(a,t)^x$.
  \end{enumerate}
\end{proof}

\begin{prop}\ollabel{prop:regular}
إذا كان~$\delta$~!!a{proof} في~\Log{N1c} أو~\Log{N1i}، وُجد
!!{proof} منتظم~$\delta'$ له البنية نفسها، وبوجه خاص~!!{height} نفسه،
ولا يختلف عن~$\delta$ إلا باستبدال المتغيرات المميزة.
\end{prop}

\begin{proof}
لأغراض البرهان وحده، سمِّ استدلال المتغير المميز في~$\delta$
\emph{نظيفًا} إذا لم يكن متغيره المميز متغيرًا مميزًا لاستدلال آخر ولم
يظهر في أي موضع من~$\delta$ خارج~!!{proof} الاستدلال الجزئي المباشر،
وسمّه \emph{متسخًا} في غير ذلك. وليكن~$n$ عدد استدلالات المتغير المميز
المتسخة في~$\delta$. نبيّن، بالاستقراء على~$n$، إمكان تحويل~$\delta$
إلى~!!{proof} المنتظم المطلوب~$\delta'$. فإذا كان~$n=0$، كانت كل
استدلالات المتغير المميز في~$\delta$ نظيفة، ومن ثم كان~$\delta$
منتظمًا ولم يبق ما نبرهنه. وإلا فاختر أعلى استدلال لمتغير مميز، وليكن
مثلًا~$\Intro{\lforall}$. يكون~!!{proof} الجزئي~$\delta_1$ من
$\delta$ المنتهي به على الصورة
    \[
    \AxiomC{}
    \RightLabel{$\delta_2$}
    \DeduceC{$!A(c)$}
    \RightLabel{\Intro{\lforall}}
    \UnaryInfC{$\lforall[x][!A(x)]$}
    \DisplayProof
    \]
لاحظ أن~$c$ لا تظهر في~$\lforall[x][!A(x)]$ ولا في فرض مفتوح لأنها
متغير مميز. ويكون البرهان~$\delta_2$ منتظمًا لأنه لا يحتوي أصلًا أي
استدلال لمتغير مميز. ولتكن~$a$~!!a{constant} لا تظهر في~$\delta$.
وبفضل~\olref{lem:subst-regular} يكون~$\Subst{\delta_2}{a}{c}$ برهانًا
منتظمًا لـ$!A(a)$. وبشرط المتغير المميز لا يحتوي أي فرض مفتوح في
$\delta_2$ الثابت~$c$، ومن ثم لا يحتوي أي فرض مفتوح في
$\Subst{\delta_2}{a}{c}$ الثابت~$a$. وهكذا يمكننا تطبيق
$\Intro{\lforall}$. وإذا استبدلنا بـ$\delta_1$ في~$\delta$ البرهان
$\delta_1'$،
    \[
    \AxiomC{}
    \RightLabel{$\Subst{\delta_2}{a}{c}$}
    \DeduceC{$!A(a)$}
    \RightLabel{\Intro{\lforall}}
    \UnaryInfC{$\lforall[x][!A(x)]$}
    \DisplayProof
    \]
نكون قد استبدلنا باستدلال متسخ لمتغير مميز استدلالًا نظيفًا، ولا فرق
إلا استبدال المتغير المميز~$c$ بـ$a$. ويكون في البرهان الناتج~$n-1$
من استدلالات المتغير المميز المتسخة، فتتبع النتيجة من فرض الاستقراء.
\end{proof}

\begin{prob}
صف الحالة في برهان~\olref{prop:regular} التي يكون فيها أعلى استدلال
لمتغير مميز هو~$\Elim{\lexists}$.
\end{prob}

لن نستند صراحة إلى القضية التي برهناها للتو، لكننا سنفترض أن كل
!!{proof}s التي نتأملها منتظمة؛ فإن لم تكن كذلك استطعنا دائمًا استبدال
متغيراتها المميزة لجعلها منتظمة.

تصدق~\Olref{lem:subst-regular} و\olref{prop:regular} أيضًا على
\Log{N2c} و\Log{N2i}. ونترك البرهانين تمرينين.

\end{document}
